% GENERATED FILE. Edit canonical lessons, then run npm run generate:books.
% canonical-source-hash: fnv1a64:f64e9523323977b1
% canonical-lessons: TE-C201-doruku, TE-C201-paduko, TE-C201-avalincu, TE-C201-vanuku, TE-C201-kadulu

\chapter{To be found, To lie down, To yawn, To shiver, To move}
\label{ch:te-to-be-found-to-lie-down-to-yawn-to-shiver-to-move}
\hlchaptermodality{201}

\begin{chapteropening}
\textbf{By the end of this chapter:} \emph{I can say to be found, to lie down, to yawn, to shiver and to move.}

Complete the last lesson of chapter 201: I can say to be found, to lie down, to yawn, to shiver and to move.
\end{chapteropening}

\section[\texorpdfstring{\te{దొరుకు} (doruku)}{దొరుకు (doruku)}]{\te{దొరుకు} (doruku) — to be found, to be available}

\label{lesson:TE-C201-doruku}



\begin{quote}
Before the new one: say the Telugu for to serve, then the Telugu for to chew.
\end{quote}

\subsection*{You'll want to know: \te{దొరుకు}}
\textbf{\te{దొరుకు}} — \emph{doruku} — \textquotedblleft{}to be found, to be available\textquotedblright{}.

\begin{grammarlens}[title={What you've built}]
One more word for this chapter.
\end{grammarlens}

\subsection*{Your turn}
\begin{itemize}
\raggedright
  \item \emph{Say it:} \emph{doruku}
  \item \emph{Say it:} \emph{doruku}, once more
  \item \emph{Say it:} say \emph{doruku}
  \item \emph{Recall:} say the Telugu for to serve, then the Telugu for to chew, then say \emph{doruku} again
\end{itemize}

\begin{tcolorbox}[breakable,colback=teal!4,colframe=teal!35!black,title={Before you move on}]
What does \te{దొరుకు} mean? (\textquotedblleft{}To be found, to be available\textquotedblright{}.) Say it once more.
\end{tcolorbox}
\section[\texorpdfstring{\te{పడుకో} (paḍukō)}{పడుకో (paḍukō)}]{\te{పడుకో} (paḍukō) — to lie down}

\label{lesson:TE-C201-paduko}



\begin{quote}
Before the new one: say the Telugu for to chew, then the Telugu for to be found.
\end{quote}

\subsection*{You'll want to know: \te{పడుకో}}
\textbf{\te{పడుకో}} — \emph{paḍukō} — \textquotedblleft{}to lie down\textquotedblright{}.

\begin{grammarlens}[title={What you've built}]
One more word for this chapter.
\end{grammarlens}

\subsection*{Your turn}
\begin{itemize}
\raggedright
  \item \emph{Say it:} \emph{paḍukō}
  \item \emph{Say it:} \emph{paḍukō}, once more
  \item \emph{Say it:} say \emph{paḍukō}
  \item \emph{Recall:} say the Telugu for to chew, then the Telugu for to be found, then say \emph{paḍukō} again
\end{itemize}

\begin{tcolorbox}[breakable,colback=teal!4,colframe=teal!35!black,title={Before you move on}]
What does \te{పడుకో} mean? (\textquotedblleft{}To lie down\textquotedblright{}.) Say it once more.
\end{tcolorbox}
\section[\texorpdfstring{\te{ఆవలించు} (āvaliñcu)}{ఆవలించు (āvaliñcu)}]{\te{ఆవలించు} (āvaliñcu) — to yawn}

\label{lesson:TE-C201-avalincu}



\begin{quote}
Before the new one: say the Telugu for to be found, then the Telugu for to lie down.
\end{quote}

\subsection*{You'll want to know: \te{ఆవలించు}}
\textbf{\te{ఆవలించు}} — \emph{āvaliñcu} — \textquotedblleft{}to yawn\textquotedblright{}.

\begin{grammarlens}[title={What you've built}]
One more word for this chapter.
\end{grammarlens}

\subsection*{Your turn}
\begin{itemize}
\raggedright
  \item \emph{Say it:} \emph{āvaliñcu}
  \item \emph{Say it:} \emph{āvaliñcu}, once more
  \item \emph{Say it:} say \emph{āvaliñcu}
  \item \emph{Recall:} say the Telugu for to be found, then the Telugu for to lie down, then say \emph{āvaliñcu} again
\end{itemize}

\begin{tcolorbox}[breakable,colback=teal!4,colframe=teal!35!black,title={Before you move on}]
What does \te{ఆవలించు} mean? (\textquotedblleft{}To yawn\textquotedblright{}.) Say it once more.
\end{tcolorbox}
\section[\texorpdfstring{\te{వణుకు} (vaṇuku)}{వణుకు (vaṇuku)}]{\te{వణుకు} (vaṇuku) — to shiver, to tremble}

\label{lesson:TE-C201-vanuku}



\begin{quote}
Before the new one: say the Telugu for to lie down, then the Telugu for to yawn.
\end{quote}

\subsection*{You'll want to know: \te{వణుకు}}
\textbf{\te{వణుకు}} — \emph{vaṇuku} — \textquotedblleft{}to shiver, to tremble\textquotedblright{}.

\begin{grammarlens}[title={What you've built}]
One more word for this chapter.
\end{grammarlens}

\subsection*{Your turn}
\begin{itemize}
\raggedright
  \item \emph{Say it:} \emph{vaṇuku}
  \item \emph{Say it:} \emph{vaṇuku}, once more
  \item \emph{Say it:} say \emph{vaṇuku}
  \item \emph{Recall:} say the Telugu for to lie down, then the Telugu for to yawn, then say \emph{vaṇuku} again
\end{itemize}

\begin{tcolorbox}[breakable,colback=teal!4,colframe=teal!35!black,title={Before you move on}]
What does \te{వణుకు} mean? (\textquotedblleft{}To shiver, to tremble\textquotedblright{}.) Say it once more.
\end{tcolorbox}
\section[\texorpdfstring{\te{కదులు} (kadulu)}{కదులు (kadulu)}]{\te{కదులు} (kadulu) — to move}

\label{lesson:TE-C201-kadulu}



\begin{quote}
Before the new one: say the Telugu for to yawn, then the Telugu for to shiver.
\end{quote}

\subsection*{You'll want to know: \te{కదులు}}
\textbf{\te{కదులు}} — \emph{kadulu} — \textquotedblleft{}to move\textquotedblright{}.

\begin{grammarlens}[title={What you've built}]
One more word for this chapter.
\end{grammarlens}

\subsection*{Your turn}
\begin{itemize}
\raggedright
  \item \emph{Say it:} \emph{kadulu}
  \item \emph{Say it:} \emph{kadulu}, once more
  \item \emph{Say it:} say \emph{kadulu}
  \item \emph{Recall:} say the Telugu for to yawn, then the Telugu for to shiver, then say \emph{kadulu} again
\end{itemize}

\begin{tcolorbox}[breakable,colback=teal!4,colframe=teal!35!black,title={Before you move on}]
What does \te{కదులు} mean? (\textquotedblleft{}To move\textquotedblright{}.) Say it once more.
\end{tcolorbox}
